\section{Preliminaries and Setting}
\label{sec:prelim}

\subsection{Words as bit-polynomials}

Fix a word width $\wbit$ (deployed: $\wbit = 32$) and the \emph{cell ring}
\[
  \cellring \;:=\; \Ftwo[X]/(X^{\wbit}),
\]
a free $\Ftwo$-module of rank $\wbit$. We identify $\cellring$ with the
bit-polynomials $\bitpoly{\wbit}$: a word $u = (u_0,\dots,u_{\wbit-1})$ is
$\bp{u} = \sum_b u_b X^b$, the coefficient of $X^b$ being bit $u_b$.
Coefficient-wise addition is bitwise $\XOR$ ($1+1 = 0$). Two further quotients
capture the operations that move bits between positions:
\[
  \Fshift = \Ftwo[X]/(X^{\wbit}) \ \text{(shift ring),}
  \qquad
  \Frot = \Ftwo[X]/(X^{\wbit}-1) \ \text{(rotation ring).}
\]
In $\Fshift$, multiplication by $X^c$ is a logical left shift with zero-fill
(high bits fall off the top); in $\Frot$ it is a cyclic rotation. The cell ring
\emph{is} the shift ring; we write $\Fshift$ when the shift structure is the
point. The three rings share the $\Ftwo$-module $\bitpoly{\wbit}$ and differ
only in how multiplication wraps.

\subsection{The binary extension field}

Fix $\K = \mathrm{GF}(2^{\kappa})$ (deployed: $\kappa = 128$, pentanomial
$X^{128}+X^7+X^2+X+1$), the protocol's \emph{challenge and evaluation field}:
all Fiat--Shamir randomness and all sumcheck arithmetic live in $\K$. Addition
is $\XOR$; multiplication is a carryless product reduced modulo the defining
polynomial. $\K$ contains $\Ftwo$ as its prime field, and we use the trivial
embedding $\Ftwo\hookrightarrow\K$ throughout.

\begin{remark}[The two roles of $X$]
\label{rem:two-X}
The symbol $X$ is the \emph{cell variable}: it indexes the $\wbit$ bit
positions of a word and, once coefficients are lifted into $\K$, the $\wbit$
coordinates of a polynomial in $\Kpoly$. It is distinct from the formal variable
defining $\K$. Every reading below substitutes the cell variable $X$ by an
element of $\K$; none touches $\K$'s internal representation.
\end{remark}

\subsection{The trace model}
\label{sec:trace}

The computation is an algebraic intermediate representation (AIR) with lookups
--- in the codebase, a \emph{universal AIR} (UAIR). A trace is a matrix of cells
$T \in \cellring^{\,N_{\mathrm{rows}}\times N_{\mathrm{cols}}}$ with
$N_{\mathrm{rows}} = 2^{\nu}$ rows (the computation's steps) and
$N_{\mathrm{cols}}$ columns. We write $\col{c}[i]$ for column $\col{c}$ at row
$i$ and $\col{c}[i\!+\!1]$ for the next row, so transition constraints relate
consecutive rows. Columns are:
\begin{itemize}[leftmargin=2em]
  \item \emph{public} (round constants, padding, the message): known to the
    verifier, not committed;
  \item \emph{primary witness}: committed by the prover;
  \item \emph{virtual}: fixed $\Ftwo$-linear functions of a primary column ---
    shifts, rotations, and $\XOR$-combinations of these (such as the
    $\Sigma/\sigma$ mixing) --- reconstructed rather than committed.
\end{itemize}
A UAIR imposes two kinds of constraint, each required at every row (transition
constraints additionally reference row $i\!+\!1$):
\begin{itemize}[leftmargin=2em]
  \item \emph{Linear / ideal-membership.} A relation
    $Q(T[i],T[i\!+\!1]) \in \ideal{g}$ asserting that an $\Ftwo[X]$-combination
    of cells lies in the ideal generated by $g$. Equality is $g = 0$; the
    modular-addition carry contributes the bit-$0$ case $g = X$; word width is
    enforced modulo $g = X^{\wbit}$. Rotations and shifts never appear here ---
    they are virtual.
  \item \emph{Hadamard.} A relation $\col{c} = \col{a}\had\col{b}$ with $\had$
    coefficient-wise multiplication of bit-polynomials --- bitwise $\AND$. This
    is the only nonlinear constraint type, and the only one that does not pass
    through the commitment linearly.
\end{itemize}
A trace \emph{satisfies} the UAIR if all constraints hold at all rows.
\Cref{sec:arith} expresses SHA-256 in exactly this language.

\subsection{Multilinear extensions}

Index rows by $\{0,1\}^{\nu}$. For a column $v\in\cellring^{\,N_{\mathrm{rows}}}$
and an $\Ftwo$-algebra $S\supseteq\cellring$, the \emph{multilinear extension}
(MLE) of $v$ over $S$ is the unique multilinear $\mle{v}\in S[Y_1,\dots,Y_\nu]$
with $\mle{v}(c) = v_c$ for $c\in\{0,1\}^\nu$:
\[
  \mle{v}(r) = \sum_{c\in\{0,1\}^{\nu}} \eq_c(r)\,v_c,
  \qquad
  \eq_c(r) = \prod_{k=1}^{\nu}\bigl(c_k r_k + (1-c_k)(1-r_k)\bigr).
\]
The PIOP works with MLEs whose coefficients are read into $\K$ and, transiently,
with MLEs valued in $\cellring$ or $\Kpoly$.

\subsection{The cell projection: one evaluation, two bases}
\label{sec:projections}

The reduction from $\Ftwo[X]$-typed constraints to a field collapses each cell
to one element of $\K$ by an $\Ftwo$-linear functional. The protocol samples
\emph{one} such functional --- evaluation at a random $\alpha$ --- but reads the
witness in two bases.

\begin{definition}[Cell projection]
\label{def:projection}
A \emph{cell projection} is an $\Ftwo$-linear map $\psi_\xi:\cellring\to\K$
determined by a weight vector $\xi = (\xi_0,\dots,\xi_{\wbit-1})\in\K^{\wbit}$
via $\psi_\xi(\sum_b w_b X^b) = \sum_b w_b\,\xi_b$. It extends $\K$-linearly to
$\Kpoly$ as the inner product of the coefficient vector with $\xi$.
\end{definition}

Both functionals the system needs are the substitution $X = \alpha$ for one
transcript-fresh $\alpha\in\K$, applied to the witness read in two bases:
\begin{itemize}[leftmargin=2em]
  \item the \emph{monomial reading}
    $\psia(w) = \sum_b w_b\alpha^b = w(\alpha)$ is genuine evaluation, an
    $\Ftwo$-\emph{algebra} homomorphism: it respects multiplication by $X$
    (shifts) and the cell ideals, and so discharges the \emph{linear}
    constraints (\Cref{sec:piop});
  \item the \emph{Lagrange reading} $\psiz(w) = \sum_b w_b L_b(\alpha)$, where
    $\{L_b\}$ is the Lagrange basis of an $\Ftwo$-linear domain $H\subset\K$ of
    size $\wbit$, turns the coefficient-wise product into a pointwise product on
    $H$, and so discharges the one nonlinear constraint, $\AND$
    (\Cref{sec:hadamard}).
\end{itemize}
These are \emph{not} two projections. There is one substitution $X = \alpha$;
$\psia$ and $\psiz$ are its values when the bits are laid into the monomial
resp.\ Lagrange basis, so even the weight vectors coincide as
$\alpha^b = X^b|_{\alpha}$ and $L_b(\alpha) = L_b(X)|_{\alpha}$. The basis is
\emph{forced by the constraint} --- monomial where multiplication-by-$X$ matters,
Lagrange where a coefficient-wise product does. Because the commitment of
\Cref{sec:commitment} binds the bit vector itself, each reading is one
length-$\wbit$ inner product against that one bound object, which is what lets
the linear and Hadamard parts share a single opening.

\begin{remark}[Why $\Ftwo$-linearity is the through-line]
\label{rem:linearity-thread}
Each $\psi_\xi$ is $\Ftwo$-linear, and the commitment's code is $\Ftwo$-linear
and does not widen $\cellring$. Hence readings commute through encoding, $\XOR$
of columns commutes with the reading, and shift/rotation virtuals commute with
both. This one fact is reused at every layer: it makes $\XOR$ free, makes shifts
virtual, lets the random reading be applied after committing, and lets one
commitment serve many functionals.
\end{remark}
